\section{Relevant State of the Art}
\label{sec:sota}

\paragraph{Naming and discovery at scale.}
The Domain Name System is the reference design for naming at Internet scale: a hierarchical name
space whose levels are separated by dots, with the data distributed so that each zone is managed by
whoever is responsible for it~\cite{rfc1034}. The instance IDs in this work reuse that idea (most specific label
first, delegation by level), but encode the \emph{physical and logical location} of an element
rather than an organisational domain. Container orchestrators such as Kubernetes provide their own
service discovery through Services and cluster DNS; in this design it does not carry the location
hierarchy to the caller, which is why a location-aware registry was kept next to the Kubernetes
Services.

\paragraph{Large-scale monitoring.}
Prometheus pulls metrics from every target and evaluates alert and recording rules locally. To
scale to many data centres its documentation describes \emph{hierarchical federation}: one
Prometheus per site scrapes local targets and a global one pulls only aggregated series from
them~\cite{prometheus_federation}. This work applies exactly this pattern, using recording rules that join
every metric with the location labels of the instance, so that the traffic between sites grows with
the number of locations rather than the number of instances.

\paragraph{SIEM.}
A SIEM combines log management, security information and event management, normalisation and
correlation of events. González-Granadillo et al.\ survey the main SIEM products and note their
convergence with big-data analytics~\cite{gonzalez2021siem}. Wazuh, the open-source SIEM used in
the practical classes, provides agents that forward logs and file-integrity events, a rule engine
with frequency-based correlation, and \emph{active responses}: scripts that run on an agent when a
rule fires, for example \code{firewall-drop}, which blocks a source with iptables for a set
time~\cite{wazuh_ar}.

\paragraph{Denial of service.}
Mirkovic and Reiher classify DDoS attacks and defences by, among others, the deployment location
of the defence (source, intermediate network, victim) and the detection strategy (pattern or
anomaly based)~\cite{mirkovic2004taxonomy}. Network-level defences include ingress filtering of
spoofed sources (BCP~38)~\cite{rfc2827} and remote-triggered blackholing of an attacked prefix
through a BGP community~\cite{rfc7999}. The Mirai botnet showed how hundreds of thousands of weak
IoT devices can generate floods that overwhelm well-defended targets~\cite{antonakakis2017mirai}.
The defence implemented here is a victim-side, per-source rate threshold; it is effective against few sources and,
as discussed later, not against a botnet.

\paragraph{Junk email.}
Content filters such as SpamAssassin combine many rules into a score; statistical (Bayesian)
filtering popularised by Graham~\cite{graham2002spam} learns from the user's own mail. Filters are
tested with GTUBE, a fixed string whose default score flags a message as spam on any
SpamAssassin installation~\cite{spamassassin_gtube}. Volume-based controls (rate limits per sender) and the refusal to
relay for unknown clients complement content filtering.

\paragraph{DNS hijacking.}
DNSSEC adds origin authentication and integrity to DNS data~\cite{rfc4033}, but does not stop an
attacker who can modify the authoritative zone itself. Such tampering happened in real campaigns:
in 2019 CISA ordered US federal agencies to audit their DNS records and monitor certificate
transparency logs after a wave of DNS infrastructure hijacking~\cite{cisa2019ed1901}. The approach taken here
targets that case (integrity of the zone at the source) with real-time file-integrity monitoring and
automatic restoration.

\paragraph{Risk analysis and fault management.}
Attack trees model the ways to reach an attacker goal as an AND/OR tree~\cite{schneier1999attacktrees};
attack--defense trees add the defender's countermeasures as nodes~\cite{kordy2010adtrees}. FCAPS
(fault, configuration, accounting, performance, security), refined by ITU-T in
M.3400~\cite{itu2000m3400}, frames fault management as detection, diagnosis, correlation and
resolution, but treats security as a separate functional area.
